\documentclass[10pt]{article}
\usepackage[margin=20mm]{geometry}
\usepackage{amsmath,amssymb,mathtools,fontspec,unicode-math,xurl,hyperref}
\IfFontExistsTF{DejaVu Serif}{\setmainfont{DejaVu Serif}}{\setmainfont{Latin Modern Roman}}
\setmathfont{latinmodern-math.otf}
\setlength{\parindent}{0pt}\setlength{\parskip}{6pt}\sloppy
\title{Technical verification of
Intermediate complex structure limit
for Calabi–Yau metrics}
\author{AI-assisted technical review}
\date{}
\begin{document}
\maketitle
Abstract: This report consolidates the preserved claim-focused review of Yang Li’s 2025 publisher version. No substantive mathematical defect is established within that review’s scope. A printed subgradient sign is corrected and checked on an admissible example. The regularization needed for the stability argument is made explicit, and a local integral estimate replaces an unsupported stronger intermediate expression. The report retains the precise dependencies of these conclusions. Scientific adjudication and independent acceptance remain pending.\par
\section*{Paper identity and review scope}
The reviewed work is Yang Li, Intermediate complex structure limit for Calabi–Yau metrics, Inventiones mathematicae 240 (2025), 459–496, DOI 10.1007/s00222-024-01314-9 (https://doi.org/10.1007/s00222-024-01314-9) [Li2025]. The source is the 38-page publisher version of record, published online on 26 February 2025. Throughout this report, PDF page means the one-based page in that file; the corresponding printed page is PDF page plus 458. Source identity and the exact SHA-256 are recorded in Appendix 7 (\#app:record).\par
This delivery continues an interrupted review. It incorporates the preserved early, middle and late audits, covering the extracted text of PDF pages 1–13, 14–26 and 27–38, respectively, and the two subsequently saved bounded checks. Those records contain the actual checks and their limits. Their work is attributed to the earlier sessions; this recovery has not repeated a full scientific review.\par
The inherited review follows the chain from the family assumptions and section decomposition, through the convex transport construction and complex ansatz, to the stability and global potential estimates. It checks normalization, signs, quantifiers and the use of cited inputs. In particular, it preserves the distinction between convergence of normalized potentials and convergence of metric tensors or microscopic geometries. The latter is not a conclusion of the main theorem.\par
\section*{Hypotheses, notation and dependencies}
In the setting of [Retained latex: \textbackslash{}cite[pp.\textasciitilde{}460--461]\{Li2025\}], let \(M\) be a smooth Fano manifold of complex dimension \(n+1\), let \(L\) be ample, and let positive integers \(d_i\) satisfy \[-K_M=\left(\sum_{i=0}^{m}d_i\right)L,\qquad 1\leq m\leq n-1.\] The sections \(F_i\) of \(d_iL\) and \(F\) of \(-K_M\) are sufficiently generic that the stated divisors are smooth and their intersections are transverse. The smooth fibers under consideration are \[X_t=\{F_0\cdots F_m+tF=0\}\subset M .\] The bound on \(m\) makes the deepest complete intersection \(E_0\cap\cdots\cap E_m\) positive-dimensional; the connectedness and section-lifting arguments use this restriction.\par
Write \[\Delta=\{x_i\geq0:\ \textstyle\sum_i x_i=1\},\qquad
 \Delta^\vee=\bigcup_{k=0}^{m}
 \{p_i\leq0,\ p_k=0,\ \textstyle\sum_i d_ip_i\geq-1\}.\] The target is understood in the quotient by \(\mathbb R(1,\ldots,1)\). The convex representation and duality in [Retained latex: \textbackslash{}cite[Eqs.\textasciitilde{}(12)--(15), p.\textasciitilde{}472]\{Li2025\}] use \[u(x)=\sup_{p\in\Delta^\vee}\{\langle x,p\rangle-\widetilde u(p)\}.\] On the target facets, the weighted measure has density \[W(p)=\left(1+\sum_i d_ip_i\right)^{n-m},\qquad
 C_1=\int_{\Delta^\vee}W(p)\,dp
 =\frac{(\sum_i d_i)(n-m)!}{n!\prod_i d_i}.\] Here \(\mu_0\) is normalized Lebesgue measure on \(\Delta\), and the intended weighted Monge–Ampère equation is \(M_u=C_1\mu_0\).\par
For the complex estimates, set \[\lambda=|\log|t||,\qquad V=(L^n),\qquad
 \omega_{\mathrm{CY},t}=\omega_t+\lambda dd^c\psi_t,\qquad
 \min_{X_t}\psi_t=0.\] The normalized Calabi–Yau measure is denoted by \(\mu_t\). The inherited normalization checks connect the weighted transport equation to the ansatz volume comparison and its total mass \(V\).\par
The review tracks the applications of the non-archimedean existence, uniqueness and domination results, Brenier transport and regularity, the uniform Skoda and stability estimates, and the Ohsawa–Takegoshi extension estimate. It does not independently reprove those external theorems. The following clarifications state exactly what the saved bounded arguments establish.\par
\section*{The subgradient sign and its checked correction}
\subsection*{Printed claim and coordinate convention}
Equation reference: eq:printed. On PDF page 16 / printed page 474, the displayed definition in Section 4.1 reads, literally, \[
 \nabla_{\rm print}u(x)=
 \{p\in\Delta^\vee:\ u(x')\geq u(x)+\langle x-x',p\rangle
       \ \hbox{for every }x'\in\Delta\}.\] This sign is inconsistent with the supporting-plane convention in (12)–(13), the Brenier map \(T=\nabla u\) on printed page 475, and the explicit formula (18) on page 476.\par
For \(m=1\), put \(x=(s,1-s)\) and \(q=p_0-p_1\). The effective coordinate \(q\) is invariant under addition of a diagonal vector. The pairing in [eq:printed] (\#eq:printed) is \((s-s')q\). At an interior differentiability point, taking \(s'=s+h\) with both signs of \(h\) forces \(q=-U'(s)\), where \(U(s)=u(s,1-s)\). The conventional subgradient has \(q=U'(s)\). A quotient-coordinate convention therefore does not remove the sign discrepancy.\par
\subsection*{Admissible example and exact discrepancy}
The saved bounded check uses \(n=2\), \(m=1\), \(d_0=1\), \(d_1=2\). These parameters are realized by a smooth quadric threefold \(M=Q^3\subset\mathbb P^4\), with \(L=\mathcal O_Q(1)\). Adjunction gives \(K_Q=(-5+2)L=-3L\). Generic sections of \(L,2L,3L\) give the required smooth transverse divisors. The resulting anticanonical pencil has smooth general fibers; shrinking the punctured parameter disk avoids its exceptional parameters. Its deepest stratum is a smooth connected curve. The example meets the stated family and dimension assumptions.\par
Equation reference: eq:example. The target interval is \(q\in[-1,1/2]\). Equation (18) of the paper gives \[
 U'(s)=
 \begin{cases}
 -1+\sqrt{3s/2},&0\leq s\leq2/3,\\[2pt]
 (1-\sqrt{3(1-s)})/2,&2/3\leq s\leq1.
 \end{cases}\] The branches are increasing, agree at zero, and have total range \([-1,1/2]\); integration yields a convex \(U\). This function also has the required representation (12): for \(q=p_0-p_1\), take \[\widetilde U(p)=p_1+\sup_{0\leq s\leq1}\{qs-U(s)\}.\] This is bounded on the compact target, and convex duality recovers \(U(s)\) as the supremum of \(\langle(s,1-s),p\rangle-\widetilde U(p)\).\par
Equation reference: eq:mass. Let \(E=(0,1/6)\). On this interval, \(-1<U'(s)<-1/2\). The literal definition [eq:printed] (\#eq:printed) requires \(q=-U'(s)\in(1/2,1)\), outside the target. Its subgradient set is empty, so its weighted image mass is zero. In contrast, \[
 C_1=\frac{1}{2d_0}+\frac{1}{2d_1}=\frac34,\qquad
 C_1\mu_0(E)=\frac34\cdot\frac16=\frac18.\] This is an exact contradiction between the literal definition and its stated application to the paper’s example.\par
\subsection*{Repair, counterevidence and impact}
Replace \(\langle x-x',p\rangle\) by \(\langle x'-x,p\rangle\). Then the image of \(E\) is \((-1,-1/2)\), on which \(W(q)=1+q\), and \[\int_{-1}^{-1/2}(1+q)\,dq=\frac18.\] On each open branch of [eq:example] (\#eq:example), \(W(U')U''=3/4\): the factors are \(\sqrt{3s/2}\) and \(3/(4\sqrt{3s/2})\) on the first branch, and \(\sqrt{3(1-s)}\) and \(3/(4\sqrt{3(1-s)})\) on the second. Continuity of \(U'\) at \(2/3\) introduces no atom there.\par
The surrounding convex-duality and transport derivation explicitly use this corrected convention. The saved result is retained as a verified printed definition/sign correction, excluded from the substantive severity tally under the task’s rule for typographical errors alone. The exact local discrepancy is preserved; it is not presented as a counterexample to the intended Calabi–Yau theorem. The reason is the absence of a demonstrated failure in the intended derivation, rather than the ease of repairing the symbol.\par
\section*{Regularization in the stability application}
\subsection*{The sufficient total-variation estimate}
The relevant passages are Proposition 5.6, Corollary 5.7 and Remark 5.8, PDF pages 23–24 / printed pages 481–482, followed by Theorems 6.1–6.3 and Corollary 6.4, pages 25–26 / 483–484. The saved check uses the local smooth approximation asserted in Remark 5.8, in addition to small potential error.\par
Fix \(t\ne0\), and write \(\mu=V\mu_t\), \(\nu=\omega_t^n\). Both measures have total mass \(V\). For a prescribed small \(\eta>0\), the good-region comparison allows \(t\) sufficiently small and a compact set \(K\) in the smooth good region such that \[\mu(K^c)\leq\eta V,\qquad
 \int_K|d\nu-d\mu|\leq\eta V .\] Passing to a compact set uses inner regularity and can be absorbed in the choice of tolerance.\par
Equation reference: eq:tv. Let \(\omega_{t,j}\) be a smooth Kähler regularization in the same class, with locally smooth convergence near \(K\) and arbitrarily small potential difference, as asserted in Remark 5.8. For the already fixed \(t,K\), choose \(j\) so that, with \(\nu_j=\omega_{t,j}^n\), \(\int_K|d\nu_j-d\nu|\leq\eta V\). Then \(a=\int_K|d\nu_j-d\mu|\leq2\eta V\), while positivity and equal total mass give \(\nu_j(K^c)\leq\mu(K^c)+a\). With the integral absolute-difference convention for total variation, \[
 \|\nu_j-\mu\|_{\mathrm{TV}}
 \leq a+\nu_j(K^c)+\mu(K^c)
 \leq2a+2\mu(K^c)\leq6\eta V .\] Thus the normalized tolerance \(\gamma^{2n+3}\) is obtained by first choosing \(\eta<\gamma^{2n+3}/6\), then \(t\), then the smoothing accuracy. No uniform smoothing scale is required.\par
\subsection*{The perturbation in Corollary 6.4}
Let \(\phi_{t,j}\) regularize the ansatz potential \(\phi_t\). For \(\omega=\omega_{\mathrm{CY},t}/\lambda\), the smooth test function is \[q_j=-\psi_t+\frac{\phi_{t,j}-\phi_t}{\lambda},
 \qquad
 \omega+dd^cq_j=\frac{\omega_{t,j}}{\lambda}.\] Choose \(\|q_j+\psi_t\|_{C^0}<\delta/8\). Then \(\sup q_j\geq-\delta/8\). If the contrary hypothesis of Corollary 6.4 holds, the set \(\{\psi_t>\delta\}\) has \(\mu_t\)-measure greater than \(\delta\), and \(q_j<-7\delta/8\) on it. The mass hypothesis of Theorem 6.3 holds with \(a=-\delta/2\), \(b=\delta\). The uniform bound on the test function is retained with a fixed margin.\par
Choose \(\gamma\) below the theorem’s threshold and so that \(C\gamma<\delta/8\). Then choose \(t\) and the regularization as above. The stability estimate implies \(\sup q_j\leq-3\delta/8\), contradicting \(\sup q_j\geq-\delta/8\). The uniform Skoda estimate, its permitted change of reference potential, and the existence of the stated regularization remain explicit analytic inputs.\par
Equation [eq:tv] (\#eq:tv) establishes the asymptotic application relative to \(V\mu_t\). It does not prove the stronger fixed-\(t\) assertion of arbitrarily small total variation relative to \(\omega_t^n\). Nor does arbitrary \(C^0\) approximation alone imply total-variation convergence. The stronger statement was not independently settled in the saved checks; no counterexample to it was established. This open verification limit does not increase supported severity.\par
\section*{The integration domain in Proposition 6.12}
On PDF page 32 / printed page 490, the displayed lower-bound chain inserts an unweighted integral over all \(X_t\) after a weight estimate known only on \(U_{y,t}\). That stronger intermediate expression is not justified by the stated local estimate. The terminal local bound needed by the proof follows directly.\par
Equation reference: eq:local. Write \(U=U_{y,t}\), \(\Phi=\phi_t+\lambda\psi_t\), and \(dV=i^{n^2}\Omega_t\wedge\overline{\Omega_t}\). The preceding local bound is \(\Phi\leq\lambda u(y)+C\lambda\varepsilon\) on \(U\). For the nonnegative weighted integrand, positivity and restriction to \(U\) give \[\begin{aligned}
 I&=\int_{X_t}|s|_{h^{\otimes l}}^2e^{-2l\Phi}\,dV
       \geq\int_U |s|_{h^{\otimes l}}^2e^{-2l\Phi}\,dV \notag\\
  &\geq e^{-2l\lambda u(y)-Cl\lambda\varepsilon}
          \int_U |s|_{h^{\otimes l}}^2\,dV .\end{aligned}\] The harmless change in \(C\) absorbs the factor two. Taking logarithms and dividing by \(2l\) yields the annular norm bound actually used in the next line of the paper. The saved late audit tracks its combination with the Ohsawa–Takegoshi upper bound and the no-cancellation estimate.\par
The finite decomposition and the point normalization \(|s(z)|=1\) select an exponent tuple using \(z\) and \(s\), independently of \(y\). This preserves the quantifier needed for the subsequent Legendre optimization. The error is \(C/l+C\varepsilon+C(l)/\lambda\), and \(t\) is chosen small after \(l\) and \(\varepsilon\) are fixed. The final optimization uses \(u(x)\geq\langle p,x\rangle-u^*(p)\), including when the actual logarithmic coordinates \(x\) do not lie on \(\Delta\).\par
The correction is to pass directly through [eq:local] (\#eq:local) and delete the unnecessary global unweighted intermediate line. It supports the already used terminal inequality under its local hypotheses. The preserved record does not establish that the stronger printed inequality is false in every instance, or that the intended core conclusion fails. This is retained as a local presentation/domain clarification, with the supported replacement shown explicitly.\par
\section*{Overall outcome and remaining scope}
The inherited checks establish no substantive Major or Minor issue within their inspected scope. The sign discrepancy and the integration-domain expression remain documented with their exact local corrections. The regularization application has the sufficient conditional argument above; its stronger fixed-\(t\) wording remains outside the independently established scope.\par
The reviewed version remains the exact publisher PDF. The saved source search identified an arXiv entry, 2305.02258, but no full version comparison is claimed. No applicable separate supplement or paper-specific released code was identified in the inspected publisher/source-link material. This is an access/search limit, not evidence of mathematical failure or a proof that no such material or erratum exists. No new empirical experiments are needed for the displayed analytic checks; no author-code execution or empirical reproduction was performed.\par
Dimension | Recorded state\par
Issue judgment | No substantive defect established in the inspected scope; two local clarifications retained.\par
Repair verification | Sign convention and explicit example checked; local integral replacement proved; stability application checked under the stated regularization and analytic inputs.\par
Unresolved scope | Stronger fixed-\(t\) regularization claim and external theorem proofs not independently established; no global correctness certification.\par
Author adoption | Not requested; no documented adoption or published correction is asserted.\par
Scientific decision | Pending. Independent acceptance remains pending. No human review is recorded.\par
1 Y. Li, Intermediate complex structure limit for Calabi–Yau metrics, Invent. Math. 240 (2025), 459–496. Publisher version of record, published online 26 February 2025. https://doi.org/10.1007/s00222-024-01314-9.\par
\section*{Supporting checks}
9a274834aa42be6427a9370a68ec0eb4
146aab312d0168b48d410ba0cc31fa57\par
The calculations in Sections 3–5 are analytic arguments reused from those records, not newly run numerical or symbolic software experiments. Their input bytes, historical origins, source locators, role attribution and current path mappings are bound in the accompanying coverage, findings and provenance files.\par
The current recovery fully parsed the 419-item input index and verified all 419 hashes, totaling 51,799,838 bytes. The read-only predecessor snapshot contains 420 files totaling 51,801,678 bytes, including its historical delivery-state file; all 420 snapshot hashes were checked. These integrity operations establish byte identity, not independent scientific acceptance.\par
The report follows the supplied 11-point, single-column `amsart` specimen. Unavailable optional font and layout packages were omitted; the existing serif mathematics and comfortable margins are retained. The command includes `–only-cached`, `–untrusted`, `–keep-logs`, `–keep-intermediates` and `–makefile-rules`. The actual complete command, exit status, source-before/source-after hashes, compiler/dependency bindings, output PDF hash and page count are recorded in the build receipt and provenance. Rendering receipts and actual page-inspection records are separate; an existing image is not treated as evidence that it was viewed.\par
\section*{Supporting checks}
Recovery-0005 supplied the two saved bounded checks and six primary source-page inspections. Those source inspections are reused only through matching source and image hashes. Recovery-0017 consolidates the saved scientific content, checks the remaining evidence-page displays, compiles the current report, and inspects its final pages.\par
These are retained as separate provenance observations, without an assertion of backend-name equivalence. No model switch is claimed. All recorded scientific reviews here are AI work. The human-review array is empty, and neither a human decision nor independent scientific acceptance is supplied by this recovery.\par
\end{document}
